\section{Verification and reproducibility}
\label{verify:section}

\subsection{Analytic proofs, finite algebra, and numerical evidence}

The proofs in the article establish the stated infinite families.
Finite symbolic tests check the coefficient algebra and the
implementation of the formulas. Numerical tests compare independent
representations and help detect phase, endpoint, or transcription
errors. A small floating-point discrepancy is not an interval
enclosure and is not used as a proof of an identity.

The complete numerical records are supplied as JSON rather than
rounded tables alone. They record working precision, truncations,
directions or endpoints, and the actual observed discrepancies.
For a reproducible source baseline, the package also includes the
archive and canonical-file hashes and a claim ledger distinguishing
new developments, inherited inputs, explicit restatements, and
unresolved questions.

\subsection{Centered harmonic and trigamma-square checks}

The file \src{code/verify_harmonic_parity.py} verifies representative
invariant and anti-invariant polynomials through inverse degree 16,
the general-shift expansions at $a=1/3$ and $a=7/5$, and the formal
trigamma shift equation through the stated truncation. It independently
checks the two finite-summation constant terms for $j=1,\ldots,8$.
These finite tests supplement the all-order injectivity proof;
they do not replace it.

At 110-digit working precision, the values $\mathcal T(-2r)$ for
$r=0,\ldots,6$ were compared with a direct finite trigamma sum plus
an Euler--Maclaurin tail. The two cutoff choices were
$(N,J)=(64,38)$ and $(96,45)$. The largest discrepancy of the
tighter-cutoff value from the theorem was approximately
$1.51\cdot10^{-85}$; the largest difference between the two
cutoff calculations was $9.28\cdot10^{-69}$.
At $t=0.2,0.7,1.1$, the polylogarithmic generator was also compared
with the Bernoulli moment series and the ordinary integral generator.
The largest moment-series discrepancy was below $4.46\cdot10^{-95}$.
No rigorous remainder enclosure is inferred from these comparisons.

\subsection{Resolvent and continuous-order checks}

The file \src{code/verify_resolvent.py} performs 55 exact assertions
on rational local coefficients, the Gamma contact polynomial, its
harmonic derivative, and the threshold where that correction
vanishes. The completed numerical run has 28 checks at 55-digit
working precision. They include upper- and lower-half-plane
parameters, twelve pointwise polygamma finite parts, the explicit
trigamma correction, the first three balanced primitives at two
parameters, two convergent $\gamma_1'$ squared-resolvent integrals,
four direct Stieltjes integrals, and two continuous polygamma orders.
A special-value regression is also included.

For the ordinary and local-subtraction checks assigned tolerance
$10^{-42}$, the largest observed discrepancy was below
$3.31\cdot10^{-53}$. The endpoint-sensitive complex order
$v=0.3+0.2i$ had discrepancy approximately $8.04\cdot10^{-41}$,
within its separately stated $10^{-34}$ tolerance.
An additional independently implemented integration-by-parts
check in \src{code/verify_resolvent_contact.py} tests nine
polygamma contact values at 100 digits, with largest discrepancy
below $5.27\cdot10^{-97}$.

The initial numerical implementation used finite differences of
the Hurwitz function at spectral order zero. At the quadrature
node $x=7/8$, that evaluation lost precision even though the
mathematical formula was exact. The final program uses the
log-Gamma identity at this order and native Hurwitz order
derivatives with an exact endpoint shift elsewhere. The regression
at $7/8$ is retained, and no tolerance was relaxed. This is an
implementation correction, not a change to the theorem.

\subsection{Harmonic elimination, Lerch continuation, and primitives}

The exact reducer \src{code/reduce_harmonic.py} is exercised by
\src{code/verify_harmonic.py}. Its recorded checks comprise
26 words, 78 comparisons among different local deletion orders
and the canonical gap form, 52 independent finite nested sums,
327 coefficient degree conditions, seven zero-insertion
identities, and 36 complete-harmonic coefficient identities.
A deliberately corrupted strict-gap sign is detected.

At 60-digit precision, six arbitrary noninteger-endpoint
continuations and one free-order derivative were evaluated by
an independent Euler--Maclaurin procedure. Their largest
observed residual was below $1.43\cdot10^{-33}$.
The zero-decoration generating function, truncated at total
degrees $10,14,18$, was compared with separate Lerch evaluations;
the residuals were approximately
$3.53\cdot10^{-20}$, $1.53\cdot10^{-26}$, and
$7.82\cdot10^{-33}$, respectively.

An ordinary Mellin integral checks the Lerch formula in its
honest convergence domain. Three normalized polynomial
primitives and five Stieltjes quadratures give additional
independent checks, with discrepancies at approximately
59--62 decimal places. A Python reciprocal initially formed
at machine precision for an exact integer input was replaced
by arbitrary-precision division; the two affected cases were
rerun and the actual corrected records retained.

\subsection{Tornheim coefficient and Mellin checks}

The file \src{code/check_quartic_tornheim.py} checks the full
quartic coefficient formula, its cyclic form, diagonal and
Euler specializations, and the elementary correction in
the mixed coordinate $\Xi$ by exact symbolic arithmetic.
It independently constructs the symmetry, axis, and
Euler-plane constraint matrices for remainder degrees
$0$ through $10$, computes their ranks and cyclic images,
and compares them with the dimension formulas.

The numerical mode uses 70-digit working precision, 68
integrated Jonqui\`ere terms, an incomplete-Gamma cutoff
of 165, and a 16-node fourth-derivative formula with step
$0.000015$. The independent Mellin comparisons were:
\begin{center}
\begin{tabular}{@{}cc@{}}
\toprule
Direction $(a,b,c)$ & Absolute discrepancy\\
\midrule
$(1,2,3)$  & $3.739835213\cdot10^{-37}$\\
$(2,3,1)$  & $2.075573689\cdot10^{-38}$\\
$(1,-2,3)$ & $2.871059279\cdot10^{-39}$\\
\bottomrule
\end{tabular}
\end{center}
Reducing the finite-difference step by a factor of ten
from an initial run decreased the first discrepancy by
the expected scale for the retained truncation error.
The delivered record is the completed tighter-step run.
The exact proof does not depend on that numerical tuning.

\subsection{How to rebuild and inspect the package}

The standalone file
\src{Harmonic_Parity_and_Resolvent_Identities.tex}
contains the complete article, including its bibliography.
It requires only a standard LaTeX installation with the
listed packages and can be compiled without network access.
The modular \src{article.tex}, \src{preamble.tex}, and
\src{sections/} files are also retained for integration.

Run \src{python3 build.py} to regenerate the standalone
source from the modules and compile the PDF.
Run \src{python3 code/verify_all.py} for the routine
exact and moderate numerical checks, or add \src{--full}
to rerun the slower independent Stieltjes, Lerch, and
Tornheim comparisons. The individual scripts remain
usable separately. Python dependencies and versions
are recorded in the package.

The PDF was compiled and inspected for unresolved references,
unavailable glyphs, overfull displays, and page layout.
The independent mathematical review notes are supplied
as review artifacts. They are an additional check by
separate derivations, not external peer review or a
machine-checked proof.
